\documentclass{article}

% Ćirilica: kodiranje fonta i jezik (srpski, ćirilica)
\usepackage[T2A]{fontenc}
\usepackage[utf8]{inputenc}
\usepackage[serbianc]{babel}

% Veličina strane i margine
\usepackage[a4paper,top=2cm,bottom=2cm,left=3cm,right=3cm,marginparwidth=1.75cm]{geometry}

% Korisni paketi
\usepackage{amsmath}
\usepackage{graphicx}
\usepackage{float}
\usepackage{booktabs}
\usepackage{tcolorbox}
\tcbuselibrary{breakable}
\usepackage[colorlinks=true, allcolors=blue]{hyperref}

% Цртежи (TikZ) и електрична кола (европски симболи)
\usepackage{tikz}
\usepackage[european]{circuitikz}
\usetikzlibrary{arrows.meta, calc, decorations.markings, babel}

% TikZ стилови који се понављају на цртежима
\tikzset{
    pomocna/.style={dashed, gray!70, thin},
    mreza/.style={gray!25, very thin},
    kriva/.style={very thick, red!75!black},
    krivaU/.style={very thick, blue!70!black},
}

% Оквир за решене примере
\newtcolorbox[auto counter]{primer}{breakable,
    colback=white, colframe=black!55, colbacktitle=black!55, coltitle=white,
    fonttitle=\bfseries\small, title={Пример~\thetcbcounter},
    boxrule=0.5pt, arc=1.5pt, left=6pt, right=6pt, top=3pt, bottom=3pt,
    before skip=8pt, after skip=10pt}

% Оквир за кључне формуле
\newtcolorbox{kljucno}{colback=blue!4, colframe=blue!45!black,
    boxrule=0.5pt, arc=1.5pt, left=8pt, right=8pt, top=4pt, bottom=4pt}

\title{ОЕТ 2 -- Отпорник, калем и кондензатор\\у колу једносмерне струје}
\author{Саобраћајно техничка школа, Београд, Земун}
\date{}

\begin{document}
\maketitle

Када у колу једносмерне струје укључимо или искључимо прекидач, струје и
напони се не морају одмах установити. Кратко време док се мењају зове се
\textbf{прелазни режим}, а стање које после тога остаје \textbf{стационарно
стање}. Отпорник реагује тренутно, а кондензатор и калем чувају енергију, па
се код њих напони и струје мењају постепено.

%==================================================================
\section{Отпорник}
%==================================================================

Отпорник енергију не чува, већ је претвара у топлоту ($P = U I = I^2 R$). Чим
затворимо прекидач, кроз њега потече струја по Омовом закону, $I = E/R$. Нема прелазног режима: струја и напон се мењају \textbf{истовремено и
скоковито}, а после тога остају стални све док је прекидач затворен.

\begin{figure}[H]
    \centering
    \begin{circuitikz}[>=Stealth, font=\small]
        \draw (0,0) to[battery1, l=$E$, invert] (0,2.6)
              to[switch, l=$\Pi$] (2.4,2.6)
              to[R, l=$R$, i>^=$i$] (2.4,0) -- (0,0);
        % --- временски дијаграм ---
        \begin{scope}[xshift=4.3cm, yshift=0.3cm]
            \draw[->] (-0.6,0) -- (5.2,0) node[right] {$t$};
            \draw[->] (0,-0.2) -- (0,2.2);
            \draw[kriva] (-0.6,0) -- (0,0) -- (0,1.6) -- (4.9,1.6);
            \node[left] at (0,1.6) {$\dfrac{E}{R}$};
            \node[below] at (0,0) {$0$};
            \node[red!75!black, above] at (2.5,1.6) {$i$};
            \node[font=\footnotesize, align=center] at (2.5,0.6) {прекидач затворен\\у тренутку $t = 0$};
        \end{scope}
    \end{circuitikz}
    \caption{Струја кроз отпорник успоставља се тренутно}
    \label{fig:R}
\end{figure}

%==================================================================
\section{Кондензатор}
%==================================================================

\subsection{Како се понаша}

Кондензатор чине две проводне плоче раздвојене изолатором (диелектриком).
Кроз изолатор струја не може да прође, али док се на плочама нагомилава
наелектрисање, у колу тече струја. Струја кроз кондензатор сразмерна је
\textbf{брзини промене} напона на њему, а напон је сразмеран наелектрисању
које је до тренутка $T$ дотекло на плоче:

\begin{kljucno}
\[
i(t) = C\,\frac{d u_C(t)}{dt},
\qquad\qquad
u_C(T) = u_C(0) + \frac{1}{C}\int_0^T i(t)\,dt
\]
\end{kljucno}

Интеграл $\int_0^T i(t)\,dt$ је наелектрисање $Q$ које је дотекло од тренутка
$0$ до тренутка $T$. Зато друга формула није ништа друго него (а), а енергија
сачувана у електричном пољу између плоча је (б):
\begin{align}
    u_C &= \frac{Q}{C} \tag{а} \\
    W   &= \frac{C\,U^2}{2} \tag{б}
\end{align}
Из прве формуле се види да \textbf{напон на кондензатору не може да се промени
скоковито}: скок напона значио би бесконачно брзу промену, а тиме и
бесконачну струју. Када се кондензатор напуни, напон се више не мења, па је
$i(t) = 0$: у стационарном стању кондензатор се понаша као \textbf{прекид}.

\textbf{Два помоћна израза} која ћемо користити (за $\tau > 0$):
\[
\frac{d}{dt}\,e^{-t/\tau} = -\frac{1}{\tau}\,e^{-t/\tau},
\qquad\qquad
\int_0^T e^{-t/\tau}\,dt = \Big[-\tau\,e^{-t/\tau}\Big]_0^T = \tau\left(1 - e^{-T/\tau}\right)
\]

\begin{figure}[H]
    \centering
    \begin{circuitikz}[font=\small]
        \draw (0,0) to[battery1, l=$E$, invert] (0,3) -- (1.4,3);
        \draw (1.4,3) node[ocirc]{} node[above]{1};
        \draw (1.4,1.9) node[ocirc]{} node[left]{2};
        \draw (1.4,1.82) -- (1.4,0) node[circ]{};
        \draw[thick] (2.3,3) -- (1.48,3.12);
        \node[above] at (1.95,3.12) {$\Pi$};
        \draw (2.3,3) node[ocirc]{} to[R, l=$R$, i>^=$i$] (5.2,3)
              to[C, l_=$C$, v^>=$u_C$] (5.2,0) -- (0,0);
    \end{circuitikz}
    \caption{Прекидач $\Pi$ у положају 1 пуни кондензатор, а у положају 2 га празни преко отпорника $R$}
    \label{fig:RC}
\end{figure}

\subsection{Пуњење}

У тренутку $t = 0$ прекидач пребацимо у положај 1. Кондензатор је празан
($u_C(0) = 0$), па је цео напон извора на отпорнику и струја је највећа: $E/R$.
Како се кондензатор пуни, $u_C(t)$ расте, напон на отпорнику опада, па опада и
струја. По другом Кирхофовом закону је

\begin{equation}
  E = R\,i(t) + u_C(t) = RC\,\frac{d u_C(t)}{dt} + u_C(t)  
\end{equation}
, а решење ове
једначине је:
\begin{kljucno}
\[
u_C(t) = E\left(1 - e^{-t/\tau}\right), \qquad
i(t) = \frac{E}{R}\,e^{-t/\tau}, \qquad
\tau = R\,C \quad [\text{s}]
\]
\end{kljucno}

\textbf{Провера изводом:} струју добијамо деривирањем напона:
\[
i(t) = C\,\frac{d}{dt}\Big[E\left(1 - e^{-t/\tau}\right)\Big]
     = C\,E\cdot\frac{1}{\tau}\,e^{-t/\tau}
     = \frac{C\,E}{R\,C}\,e^{-t/\tau}
     = \frac{E}{R}\,e^{-t/\tau}
\]

За $T \to \infty$ на кондензатор дотекне наелектрисање
$Q = \int_0^\infty i(t)\,dt = \frac{E}{R}\,\tau = C\,E$, као што и треба.

$\tau$ је \textbf{временска константа}. Што је већа, пуњење траје дуже. После
$\tau$ кондензатор је напуњен до $63\,\%$ од $E$, после $3\tau$ до $95\,\%$, а
после $5\tau$ до $99\,\%$, па у пракси сматрамо да је напуњен после $5\tau$.

\subsection{Пражњење}

Када прекидач пребацимо у положај 2, напуњени кондензатор ($u_C(0) = E$)
постаје извор и празни се преко отпорника. Сада у контури нема извора, па је
$0 = R\,i(t) + u_C(t)$. Струја тече у \textbf{супротном смеру} од струје
пуњења, па је негативна (време $t$ меримо од тренутка пребацивања прекидача):
\begin{kljucno}
\[
u_C(t) = E\,e^{-t/\tau}, \qquad i(t) = -\frac{E}{R}\,e^{-t/\tau}
\]
\end{kljucno}

\textbf{Провера изводом:}
\[
i(t) = C\,\frac{d}{dt}\Big[E\,e^{-t/\tau}\Big]
     = C\,E\cdot\left(-\frac{1}{\tau}\right) e^{-t/\tau}
     = -\frac{E}{R}\,e^{-t/\tau}
\]

\begin{figure}[H]
    \centering
    \begin{tikzpicture}[>=Stealth, font=\small]
        \def\k{1.15}  % cm по τ
        \def\U{1.6}   % E на цртежу
        \def\I{0.9}   % E/R на цртежу
        % осе
        \draw[->] (-0.4,0) -- ({10*\k+0.5},0) node[right] {$t$};
        \draw[->] (0,-1.3) -- (0,2.1);
        \draw[pomocna] (0,\U) -- ({10*\k},\U);
        \node[left] at (0,\U) {$E$};
        \node[left] at (0,\I) {$\frac{E}{R}$};
        \node[left] at (0,-\I) {$-\frac{E}{R}$};
        \draw[pomocna] (0,-\I) -- (0.3,-\I);
        % граница пуњење / пражњење
        \draw[dashed] ({5*\k},-1.25) -- ({5*\k},2.05);
        \node[font=\footnotesize] at ({2.5*\k},2.25) {пуњење ($\Pi$ у 1)};
        \node[font=\footnotesize] at ({7.5*\k},2.25) {пражњење ($\Pi$ у 2)};
        % τ
        \draw[pomocna] (\k,0) -- (\k,{0.632*\U}) -- (0,{0.632*\U});
        \node[below, font=\footnotesize] at (\k,0) {$\tau$};
        \node[above left, font=\footnotesize, inner sep=1pt] at (\k,{0.632*\U}) {$0{,}63E$};
        \node[below right, font=\footnotesize] at ({5*\k},0) {$5\tau$};
        % напон u_C
        \draw[krivaU, domain=0:5, samples=60] plot ({\x*\k},{\U*(1-exp(-\x))});
        \draw[krivaU, domain=5:10, samples=60] plot ({\x*\k},{\U*(1-exp(-5))*exp(-(\x-5))});
        \node[blue!70!black, above] at ({3.4*\k},\U) {$u_C$};
        % струја i
        \draw[kriva] (-0.4,0) -- (0,0) -- (0,\I);
        \draw[kriva, domain=0:5, samples=60] plot ({\x*\k},{\I*exp(-\x)});
        \draw[kriva] ({5*\k},0) -- ({5*\k},-\I);
        \draw[kriva, domain=5:10, samples=60] plot ({\x*\k},{-\I*exp(-(\x-5))});
        \node[red!75!black, above right] at ({1.4*\k},{\I*exp(-1.4)}) {$i$};
    \end{tikzpicture}
    \caption{Напон $u_C$ и струја $i$ при пуњењу и пражњењу кондензатора}
    \label{fig:RC-dijagram}
\end{figure}

\begin{primer}
Кондензатор $C = 100\,\mu$F пуни се преко отпорника $R = 10$~k$\Omega$ из
извора $E = 12$~V. Колика је временска константа, колика је почетна струја и
за које време се кондензатор практично напуни?

\textbf{Решење.}
$\tau = R\,C = 10^{4}~\Omega \cdot 10^{-4}~\text{F} = 1$~s, \
$i(0) = \frac{E}{R} = \frac{12~\text{V}}{10^4~\Omega} = 1{,}2$~mA.
Кондензатор је практично напуњен после $5\tau = 5$~s, а после $1$~s напон на
њему је $0{,}63\cdot 12 \approx 7{,}6$~V.
\end{primer}

%==================================================================
\section{Калем}
%==================================================================

\subsection{Како се понаша}

Калем је намотана жица. Струја кроз њега ствара магнетно поље. Када се
струја мења, мења се и поље, па се у калему индукује напон који се
\textbf{супротставља промени струје}. Напон на калему сразмеран је брзини
промене струје, а струја је сразмерна интегралу напона од $0$ до $T$:

\begin{kljucno}
\[
u_L(t) = L\,\frac{d i(t)}{dt},
\qquad\qquad
i(T) = i(0) + \frac{1}{L}\int_0^T u_L(t)\,dt
\]
\end{kljucno}

$L$ је \textbf{индуктивност} (јединица хенри, H). Интеграл
$\int_0^T u_L(t)\,dt$ је промена магнетног флукса $N\Phi$ који обухвата
намотај (по Фарадејевом закону $u_L(t) = N\,\frac{d\Phi(t)}{dt}$). Зато друга
формула није ништа друго него (а), а енергија сачувана у магнетном пољу је (б):
\begin{align}
    L &= \frac{N\,\Phi}{I} \tag{а} \\
    W &= \frac{L\,I^2}{2} \tag{б}
\end{align}

Из прве формуле се види да \textbf{струја кроз калем не може да се промени
скоковито}: скок струје значио би бесконачан напон. Када се струја устали,
$\frac{d i(t)}{dt} = 0$, па је и $u_L(t) = 0$: у стационарном стању калем се
понаша као \textbf{кратак спој} (остаје само мала отпорност жице).

Калем и кондензатор су „супротни”: кондензатор чува напон, а калем струју.

\subsection{Укључење}

Коло је исто као на слици~\ref{fig:RC}, само уместо кондензатора стоји калем
$L$. Када прекидач пребацимо у положај 1, струја не може одмах да скочи, па
у првом тренутку ($i(0) = 0$) цео напон извора пада на калему. Струја
постепено расте, а напон на калему опада. По другом Кирхофовом закону је

\begin{equation}
    E = R\,i(t) + L\,\frac{d i(t)}{dt}
\end{equation}

, а решење је:
\begin{kljucno}
\[
i(t) = \frac{E}{R}\left(1 - e^{-t/\tau}\right), \qquad
u_L(t) = E\,e^{-t/\tau}, \qquad
\tau = \frac{L}{R} \quad [\text{s}]
\]
\end{kljucno}

\textbf{Провера изводом:}
\[
u_L(t) = L\,\frac{d}{dt}\Big[\frac{E}{R}\left(1 - e^{-t/\tau}\right)\Big]
       = L\,\frac{E}{R}\cdot\frac{1}{\tau}\,e^{-t/\tau}
       = L\,\frac{E}{R}\cdot\frac{R}{L}\,e^{-t/\tau}
       = E\,e^{-t/\tau}
\]

\subsection{Искључење}

Када прекидач пребацимо у положај 2, калем одржава струју у истом смеру
($i(0) = E/R$) и она постепено опада кроз отпорник. У контури нема извора, па
је $0 = R\,i(t) + L\,\frac{d i(t)}{dt}$. Калем се сада понаша као извор, па
напон на њему мења знак:
\begin{kljucno}
\[
i(t) = \frac{E}{R}\,e^{-t/\tau}, \qquad u_L(t) = -E\,e^{-t/\tau}
\]
\end{kljucno}

\textbf{Провера изводом:}
\[
u_L(t) = L\,\frac{d}{dt}\Big[\frac{E}{R}\,e^{-t/\tau}\Big]
       = L\,\frac{E}{R}\cdot\left(-\frac{R}{L}\right) e^{-t/\tau}
       = -E\,e^{-t/\tau}
\]

Ако коло нагло прекинемо, а струја нема куда да тече, она пада готово
тренутно. Тада је промена струје огромна, па се на калему индукује
\textbf{врло висок напон} и између контаката прекидача скочи варница. Ова
појава се користи у \textbf{паљењу мотора}: бобина (калем) при прекиду струје
даје напон од више хиљада волти за свећицу.

\begin{figure}[H]
    \centering
    \begin{tikzpicture}[>=Stealth, font=\small]
        \def\k{1.15}  % cm по τ
        \def\U{1.6}   % E на цртежу
        \def\I{0.9}   % E/R на цртежу
        % осе
        \draw[->] (-0.4,0) -- ({10*\k+0.5},0) node[right] {$t$};
        \draw[->] (0,-1.8) -- (0,2.1);
        \draw[pomocna] (0,\I) -- ({10*\k},\I);
        \node[left] at (0,\U) {$E$};
        \node[left] at (0,-\U) {$-E$};
        \node[left] at (0,\I) {$\frac{E}{R}$};
        % граница укључење / искључење
        \draw[dashed] ({5*\k},-1.75) -- ({5*\k},2.05);
        \node[font=\footnotesize] at ({2.5*\k},2.25) {укључење ($\Pi$ у 1)};
        \node[font=\footnotesize] at ({7.5*\k},2.25) {искључење ($\Pi$ у 2)};
        % τ
        \draw[pomocna] (\k,0) -- (\k,{0.632*\I});
        \node[below, font=\footnotesize] at (\k,0) {$\tau$};
        \node[below right, font=\footnotesize] at ({5*\k},0) {$5\tau$};
        % струја i
        \draw[kriva, domain=0:5, samples=60] plot ({\x*\k},{\I*(1-exp(-\x))});
        \draw[kriva, domain=5:10, samples=60] plot ({\x*\k},{\I*exp(-(\x-5))});
        \node[red!75!black, above] at ({3.4*\k},\I) {$i$};
        % напон u_L
        \draw[krivaU] (-0.4,0) -- (0,0) -- (0,\U);
        \draw[krivaU, domain=0:5, samples=60] plot ({\x*\k},{\U*exp(-\x)});
        \draw[krivaU] ({5*\k},0) -- ({5*\k},-\U);
        \draw[krivaU, domain=5:10, samples=60] plot ({\x*\k},{-\U*exp(-(\x-5))});
        \node[blue!70!black, above right] at ({1.3*\k},{\U*exp(-1.3)}) {$u_L$};
    \end{tikzpicture}
    \caption{Струја $i$ и напон $u_L$ при укључењу и искључењу калема}
    \label{fig:RL-dijagram}
\end{figure}

\begin{primer}
Калем $L = 0{,}5$~H везан је редно са отпорником $R = 10\,\Omega$ на извор
$E = 12$~V. Колика је временска константа, колика је струја у стационарном
стању и колика је тада енергија у калему?

\textbf{Решење.}
$\tau = \frac{L}{R} = \frac{0{,}5}{10} = 0{,}05$~s $= 50$~ms, \
$I = \frac{E}{R} = 1{,}2$~A, \
$W = \frac{L\,I^2}{2} = \frac{0{,}5\cdot 1{,}2^2}{2} = 0{,}36$~J.
Струја се практично устали после $5\tau = 250$~ms.
\end{primer}

%==================================================================
\section{Поређење}
%==================================================================

\begin{table}[H]
    \centering
    \small
    \renewcommand{\arraystretch}{1.5}
    \begin{tabular}{lccc}
        \toprule
                                   & \textbf{Отпорник} & \textbf{Кондензатор} & \textbf{Калем} \\
        \midrule
        не може скоковито да се промени & --          & напон $u_C$          & струја $i$ \\
        у тренутку укључења        & $i = E/R$ одмах   & као кратак спој      & као прекид \\
        у стационарном стању       & $i = E/R$         & прекид ($i = 0$)     & кратак спој \\
        временска константа        & --                & $\tau = R\,C$        & $\tau = L/R$ \\
        енергија                   & топлота $I^2 R\,t$ & $\dfrac{C U^2}{2}$ (ел. поље) & $\dfrac{L I^2}{2}$ (маг. поље) \\
        \bottomrule
    \end{tabular}
    \caption{Понашање елемената у колу једносмерне струје}
    \label{tab:poredjenje-dc}
\end{table}

\end{document}